\chapter{Discussão e Conclusão}\label{ch:discussion}

\section{Interpretação dos resultados}\label{sec:interpretation}

O zoneamento responde à QP1: o agrupamento descorrelacionado resolve GO~e~DF em dez zonas
biofisicamente distintas, que se agrupam em seis vocações comparativas --- um planalto alto e
argiloso voltado à pecuária (Zona~1); um planalto intermediário voltado a outras culturas anuais
(Zona~3); o núcleo agrícola de soja e cana-de-açúcar (Zona~10); duas planícies arenosas voltadas
à geração fotovoltaica (Zonas~6 e~9, a segunda também a mais árida do território); duas zonas
voltadas à piscicultura, uma baixada ripária (Zona~2) e uma baixada quente e remota (Zona~7); e
três zonas de conservação de perfis físicos bem distintos entre si --- um vale úmido de
transição (Zona~4), um planalto periurbano (Zona~8) e um planalto alto e íngreme onde a
viabilidade de culturas colapsa e a conservação é também o melhor uso absoluto (Zona~5, com
74\% de Cerrado nativo). A soja, apesar de ter viabilidade absoluta alta em quase toda zona de
cultivo, nunca é o melhor uso comparativo de nenhuma delas. Essa partição usa apenas
atributos biofísicos, nunca a cobertura da terra, e ainda assim reproduz a geografia
observada do uso do território, o que sustenta sua plausibilidade. O padrão subjacente é que
GO~e~DF é climaticamente quase uniforme: quem separa terra boa de terra ruim é o solo e o
terreno, não o clima. Essa leitura carrega, portanto, o peso de dois produtos globais herdados
sem validação local: o solo vem do OpenLandMap a 250~m, e o terreno do SRTM GL1 bruto a 30~m
\citep{Farr2007}, reamostrado para a mesma grade; é a acurácia desse par de camadas, não a do
clima, que se propaga diretamente para as fronteiras finas entre zonas (detalhado
em~\S\ref{sec:limitations}).

A projeção climática estende essa leitura a meados do século (QP3). A mudança é de magnitude
pequena a moderada, mas direcionalmente coerente em todo o conjunto de modelos, e para as
culturas ela aguça o gradiente já existente em vez de redesenhá-lo: as perdas recaem sobre a
baixada quente, arenosa e remota já marginal (Zona~7), enquanto os planaltos privilegiados
permanecem essencialmente inalterados. Três segmentos são mais expostos. A conservação é o mais
afetado: sua área viável cai de 27,8\% para 12,8\% no cenário--janela mais quente, um resultado
coerente com o fato de a vegetação nativa ser o uso da terra menos amortecido pelo clima, ao
contrário das culturas, mantidas perto de seu ótimo térmico pelo manejo. As outras culturas
anuais (milho, algodão, sorgo) recuam de 58,3\% para 51,0\%, e a geração fotovoltaica, de 20,5\%
para 14,4\% por perda de desempenho térmico do painel. Soja e cana-de-açúcar permanecem
praticamente intactas.

A fase de uso atual da terra responde à QP2 e, de quebra, valida a QP1: a soja atual alcança 32\% do planalto de argila alta (Zona~1)
contra parcelas bem menores nas zonas marginais, e as culturas ocupam terra sistematicamente
mais viável que a média disponível (soja: 0,935 nos pixels plantados contra 0,773 na terra
apta). A discriminação por presença é forte para a soja sobre uma amostra estratificada (AUC~0,871, Boyce~0,964).

Duas tensões merecem registro. Primeira, uma lacuna substancial separa potencial e uso atual:
41,8\% do território é viável para soja e ainda não cultivado, majoritariamente sobre pastagem.
Isso aponta para intensificação, não para abertura de novas áreas, o que é coerente com a
trajetória já documentada para a fronteira agrícola do Cerrado \citep{Spera2017,Caballero2023}.
Segunda, a verificação cruzada orientada por dados encontra o mapa baseado em conhecimento mais
permissivo que o nicho realizado ($\kappa=+0{,}269$, concordância razoável mas não forte na
escala de Landis \& Koch; o mapa de conhecimento marca 52,8\% do território como viável,
contra 21,6\% do classificador orientado por dados). Como o clima é quase uniforme, os
limiares fuzzy não podem ser estreitados por ele; parte dessa permissividade residual é
provavelmente regra especialista ainda um pouco larga, e parte é o próprio sinal da QP2: terra
capaz deixada sem uso por razões econômicas, de posse ou de logística que o modelo não capta.

Por que esse $\kappa$ fica abaixo do que a AUC de presença isoladamente sugeriria (0,871) é
explicado em \S\ref{sec:validation-present}: são, no fundo, duas perguntas diferentes sobre o
mesmo resultado.

Essa leitura econômica tem respaldo na literatura sobre a fronteira agrícola do Cerrado. Em
Goiás, a intensificação registrada entre 2000 e 2016, e não a abertura de novas áreas,
respondeu por boa parte do ganho de produção: a lotação do rebanho bovino subiu de 1,17 para 1,24
cabeças/hectare e a área de dupla safra soja-milho quase quadruplicou no período, ainda que entre
50\% e 60\% das pastagens do Cerrado sigam degradadas \citep{Spera2017}. Recuperar essa pastagem
para lavoura ou pecuária mais intensiva tem um custo de entrada que a viabilidade biofísica, por
si só, não precifica --- calagem e correção da acidez do solo, mecanização, acesso a crédito;
\citet{Spera2017} reporta ainda que cerca de 40\% dos produtores brasileiros não têm título de posse da
terra, o que restringe o acesso a crédito e, por essa via, a capacidade de intensificar mesmo onde
o solo permite. Em escala de bioma, a conversão de pastagem para soja e outras culturas responde
também a condições de mercado --- preços de commodities, câmbio, demanda por alimentos e
biocombustíveis --- mais do que à mera disponibilidade de terra apta \citep{Caballero2023}. Preço,
custo de recuperação do solo, posse da terra e logística de escoamento são, portanto, candidatos
plausíveis para explicar a lacuna de 41,8\% aqui observada, mas nenhum deles é modelado neste
trabalho; testá-los é o objeto da camada econômica proposta em \S\ref{sec:futurework}.

\section{Contribuições metodológicas}\label{sec:contributions}

Além dos mapas em si, algumas escolhas de desenho valem registro. O procedimento roda
inteiramente na nuvem, a partir do catálogo público e de uma especificação de parâmetros sob
controle de versão, sem custo de infraestrutura dedicada; isso o torna reexecutável para outros
estados do Cerrado. Os sete segmentos são avaliados por um mecanismo comum de pertinência
fuzzy/AHP e comparados em base normalizada (escore~z), com a regra de agregação ajustada a cada
caso: média geométrica (fator limitante) para os segmentos de produção, e média aritmética
ponderada para a conservação, tratada como índice de valor e não de estado, construída a
partir de critérios de conectividade, heterogeneidade de habitat e serviço de carbono/água, em
vez de simplesmente repetir a cobertura nativa atual.

O agrupamento k-means sobre o conjunto bruto de 34 bandas colapsava em poucas zonas pouco
informativas, dominado pelas bandas climáticas colineares; reduzir, ponderar por tema e aplicar
PCA antes do agrupamento \citep{Ding2004} resolveu as dez zonas relatadas. E reutilizar a mesma
transformação de derivação climática tanto para a linha de base quanto para o futuro do CMIP6
garante que a mudança projetada seja um sinal puramente climático (fatores de mudança identidade
dão $\Delta S=0$), com os demais fatores mantidos estáticos para isolar essa componente. A
comparação com floresta aleatória sob validação cruzada espacial (\S\ref{sec:interpretation})
serve de checagem independente do mapa baseado em conhecimento, não de substituto: as
divergências encontradas são tratadas como achado sobre potencial \emph{versus} realizado, não
como defeito oculto.

\section{Limitações}\label{sec:limitations}

Parte das limitações abaixo decorre da escolha deliberada pelo desenho descrito na seção anterior; 
outras são lacunas genuínas para trabalho futuro.

\begin{itemize}
  \item \textbf{Escopo biofísico, não econômico.} O modelo capta potencial agroclimático mais
    um \emph{proxy} de acessibilidade por tempo de viagem, não preços, posse da terra,
    crédito ou logística. Parte da lacuna potencial \emph{versus} realizado é exatamente essa
    geografia econômica ausente.
  \item \textbf{Solo e terreno de produtos globais únicos.} A decomposição por fator mostra que,
    com o clima quase uniforme, solo e terreno carregam a maior parte da discriminação espacial
    (\S\ref{sec:interpretation}); o atlas herda, portanto, a acurácia de dois produtos globais que não
    possuem um ajuste fino a realidade local. O solo vem do
    OpenLandMap a 250~m. O terreno vem do SRTM GL1 bruto a 30~m \citep{Farr2007}, não uma
    versão com preenchimento de vazio como o NASADEM ou o CGIAR-CSI v4.1, com acurácia
    vertical absoluta especificada em~16~m (\emph{linear error} a 90\% de confiança) e vazios de
    sombra de radar/água que o \emph{pipeline} não trata antes de reamostrar para 250~m por média
    (elevação, declive, orientação) ou interpolação bilinear (índice topográfico de umidade).
 
  \item \textbf{Descompasso de escala clima--solo.} As entradas climáticas ($\approx$4,6~km,
    TerraClimate) são bem mais grosseiras que solo e terreno (250~m). Um levantamento do
    catálogo do GEE não encontrou alternativa viável: WorldClim é mais fino ($\approx$927~m) mas
    sua climatologia (1960--1990) não cobre a normal 1991--2020 usada aqui e não oferece
    evapotranspiração potencial (PET), déficit de pressão de vapor (VPD) ou umidade do solo;
    CHIRPS/CHIRTS cobrem uma única variável cada e, a 0,05\textdegree{} ($\approx$5,6~km), são
    mais grosseiros que o TerraClimate; ERA5-Land/AgERA5 têm o conjunto completo de variáveis e
    o período correto, mas ficam a 0,1\textdegree{} ($\approx$9--11~km, mais de 2$\times$ pior),
    e o AgERA5 sequer está no catálogo oficial. Esse descompasso nominal de resolução é
    enganoso: o clima atua como comporta, não como fonte de estrutura fina, já demonstrado
    quantitativamente pelo mapa de magnitude da Figura~\ref{fig:scale-mismatch}
    (\S\ref{sec:atlas-present}). Essa leitura, porém, simplifica demais se aplicada por
    \emph{prefixo de origem do dado} em vez de por \emph{papel temático}: na decomposição por
    fator agregada por tema (Figura~\ref{fig:factor-variance}), o clima explica em média 0,13 da
    variância territorial, mas essa média esconde dois extremos --- a irradiância solar
    (\texttt{clim\_srad}, 0,36) é um sinal climático espacial genuíno, enquanto o fator
    nominalmente ``climático'' que domina a pecuária (\texttt{clim\_soil\_moist}, 0,58) descreve
    armazenamento de água no solo, não um gradiente climático; reclassificado pelo papel
    temático, a parcela de terreno+solo da pecuária sobe de 0,42 para 1,00. A nomenclatura de
    origem do dado não deve, portanto, ser confundida com o papel temático da variável: é a
    padronização (\S\ref{sec:fuzzy-standard}) e a agregação por fator que definem o que conta
    como ``clima'', não o catálogo de onde a banda vem.
  \item \textbf{Sobrepermissividade residual das culturas.} As superfícies de culturas
    permanecem mais permissivas que o nicho realizado (\S\ref{sec:interpretation},
    $\kappa=+0{,}269$ para a soja). Separar regra especialista larga de lacuna potencial
    \emph{versus} realizado exigiria, ainda assim, um alvo de calibração ancorado em rendimento,
    que este desenho não usa.
  \item \textbf{Validador de produtividade é um substituto.} O NPP/GPP do MODIS mede produção
    primária da vegetação, não rendimento agronômico \citep{Turner2006,He2018}, e a confusão
    entre coberturas limita seu uso a contexto descritivo; AUC e Boyce seguem como validadores
    primários de presença de cultura.
  \item \textbf{Piscicultura tem uso realizado quase inexistente no MapBiomas.} Apenas um 
    único pixel de aquicultura é mapeado pelo MapBiomas em todo o território GO~e~DF, 
    amostra insuficiente para ancorar limiares ou validar presença contra esse uso realizado. 
    A calibração da piscicultura permanece ancorada na terra genericamente \emph{disponível}, 
    não no nicho realizado, e nenhum validador de presença (Boyce/AUC) é
    reportado para esse segmento pela mesma razão. Resolver esta lacuna envolveria obter uma
    nova classificação da área utilizada para psicultura na área de interesse.
  \item \textbf{Projeção futura.} A projeção CMIP6, como qualquer
    projeção por fatores de mudança \citep{RamirezVillegas2010,Hay2000}, só pode ser sustentada
    indiretamente (linha de base validada, concordância entre GCMs); radiação solar e umidade do
    solo foram mantidas no valor histórico por falta de um campo projetado e de um balanço
    hídrico completo, o que atenua o sinal climático da geração fotovoltaica e da pecuária.
\end{itemize}

\section{Trabalhos futuros}\label{sec:futurework}

%% TODO citar IBGE PAM e PPM como fontes para validadores economicos reais, possivel estimação de perda economica
Uma primeira frente é substituir as duas variáveis climáticas mantidas estáticas, introduzindo
radiação solar projetada e um tratamento de umidade do solo por balanço hídrico, para que
geração fotovoltaica e pecuária passem a carregar sinal climático completo. Outra é reforçar a
validação: incorporar um validador ancorado em rendimento (estatísticas agrícolas municipais, por
exemplo) e estender a verificação cruzada com floresta aleatória e blocos espaciais a todos os
segmentos.
Uma terceira direção é acrescentar uma camada econômica (preços, posse da terra, infraestrutura),
que converteria parte da lacuna potencial \emph{versus} realizado de resíduo em quantidade modelada.
Uma quarta é explorar modelos espacialmente explícitos (por exemplo, regressão com defasagem
espacial ou um termo autorregressivo condicional) para capturar efeitos de vizinhança entre
células: o mecanismo fuzzy/AHP atual trata cada pixel de forma independente, e já produz uma
interpretação espacial coerente (\S\ref{sec:interpretation}), mas não modela explicitamente a
dependência espacial entre células vizinhas, que poderia refinar tanto o zoneamento quanto a
leitura de incerteza da projeção futura.

Por fim, cabe revisar os critérios de
valor de conservação e a direção de sua resposta climática (perda de viabilidade \emph{versus}
prioridade crescente de proteção), no espírito do planejamento de compromissos entre conservação e
restauração já aplicado ao Cerrado \citep{Schuler2022}.

\section{Conclusão}\label{sec:conclusion}

Esta dissertação mapeou o potencial agroclimático de sete segmentos de uso da terra em Goiás e
no Distrito Federal, particionou o território em zonas agroambientais e projetou como esse
potencial se desloca sob mudança do clima até meados do século, com um método aberto e
reprodutível na nuvem.

As três questões de pesquisa foram respondidas. Um mecanismo fuzzy/AHP produz superfícies de
viabilidade do presente para os sete segmentos, e um agrupamento não supervisionado resolve o
território em dez zonas que reproduzem, de forma independente, a geografia observada do uso da
terra. A projeção CMIP6 mostra uma mudança modesta mas direcionalmente robusta até 2070, com a
conservação como o segmento mais exposto, seguida das outras culturas anuais e da geração
fotovoltaica, enquanto soja e cana-de-açúcar permanecem estáveis (QP3). Uma lacuna espacialmente
validada entre potencial e uso atual, concentrada em pastagem existente, aponta para
intensificação como direção mais provável de mudança de uso da terra na região.

De forma mais geral, o trabalho mostra que uma avaliação de terras multissegmento, sensível ao
clima e validada contra o uso atual pode ser montada inteiramente a partir de dados públicos, na
nuvem, a custo marginal baixo. O atlas resultante indica onde cada segmento tem melhor aptidão em
GO~e~DF, onde esse potencial ainda não é realizado e como a mudança do clima deve redesenhá-lo até
meados do século.
